\subsection{Implementation Specifics of Corporate Website Scenario}\label{subsec:implementation-specifics-of-corporate-website-scenario}
After setting up the framework most of the features are already implemented and incorporated by the initial download.
The only thing left to implement is the homepage.

Since there should only be one instance of homepage in the system, the framework does not provide separate template for
the homepage and all the implementation details are contained in the src/components/Home directory.
The directory includes Home.php file that describes the interface of the component with private functions used only
within the Home.phtml template.
When the template is rendered, it imports the home.css styles and then renders the HTML of the page.

The Home directory provided by the framework does not include JavaScript file.
The homepage requires JavaScript functionality and thus home.js is created and edited.
The template is edited to import the script file using SideLoader's Javascript::import() method and the homepage layout
is implemented.

\subsubsection*{Testing section: CMS Components}
All the CMS components needed for the website are already implemented and functional to the extent the application needs.
The only exception is manual edit of homepage which is encouraged by the framework's implementation.

\subsubsection*{Testing section: AI Generation}
The framework provides couple of ways to generate content more efficiently by using AI\@.
The following table describes whether the feature was used and if the output matched the expectation.

\begin{table}[H]
    \centering
    \begin{tabular}{|p{0.4\linewidth}|p{0.1\linewidth}|p{0.4\linewidth}|}
        \hline
        \textbf{Feature}
        &
        \textbf{Used}
        &
        \textbf{Output quality} \\
        \hline
        Page Structure Generation
        &
        Yes
        &
        The generated structure is usable but a few pages need to be removed.
        For instance Home needs to be removed because it is implemented via the Home component. \\
        \hline
        Article Generation
        &
        Yes
        &
        The generated output matched the desired content but the Markdown formatting had to be manually edited
        because the framework uses custom-made Markdown compiler and some the formatting choices lead to undefined
        results. \\
        \hline
        Static Page Generation
        &
        Yes
        &
        The static contact page was generated using AI with simple prompt that does not include any styling guides
        and detailed description of the website as a whole.
        The effect of low quality prompt and smaller model leads to underwhelming results. \\
        \hline
        Page Component Generation
        &
        No
        & \\
        \hline
        Documentation Generation
        &
        No
        &
        \\
        \hline
    \end{tabular}
    \caption{Personal Testing Section - AI Generation}
    \label{tab:startuh-ai}
\end{table}

\subsubsection*{Testing section: Metrics Evaluation}
The framework is tested on the metrics defined in the table.
Some values are missing metrics about bug fixing the framework.
If critical bug is encountered during the deployment, the development of fixing the bug is not counted to the metrics.

\begin{table}[H]
    \centering
    \begin{tabular}{|p{0.4\linewidth}|p{0.5\linewidth}|}
        \hline
        \textbf{Evaluation Metric}
        &
        \textbf{Value} \\
        \hline
        Requirement Coverage
        &
        The only thing that had to be implemented by hand is the the homepage and since it is intended, the Requirement
        Coverage is 100\% \\
        \hline
        Number of Manual Interventions
        &
        Only one, homepage implementation \\
        \hline
        Time to Publishable Draft
        &
        Measuring the time between initial commit and final commit (not counting update commits) the deployment of the website
        took exactly one hour and two minutes. \\
        \hline
        Responsiveness
        &
        All the pages that are visible to the public are fully responsive although some elements are implemented
        in functional but not user-friendly way.
        The elements are header menu items on mobile devices.
        The administration websites require custom implementation of core user controls, and thus it is only optimized
        for desktop devices with standard landscape resolutions.
        \\
        \hline
        Basic SEO
        &
        Publicly accessible websites pass the SEO test if all metadata for a page are filled in.
        The only exception is the homepage which requires manual description to be added by editing the Home component
        directly.
        \\
        \hline
    \end{tabular}
    \caption{Personal Testing Section - Metrics Evaluation}
    \label{tab:startuh-evaluation}
\end{table}
